
\subsubsection{Summary of Findings}

The mechanism chain (annotator conflation $\rightarrow$ reward conflation $\rightarrow$ representation conflation) is verified through H-M1, H-M2, and H-M3. The H-M4 failure reflects detection inadequacy rather than mechanism failure---keyword-based features captured only 2.3\% of tasks, providing insufficient statistical power to test the correlation.

\subsubsection{Interpretation}

The experiments establish that RLHF training creates a reward model conflation pattern:

\begin{enumerate}
\item Human annotators rate both correctness and user-state-modeling tasks with similar high-confidence patterns (rate\_diff = 0.001).
\end{enumerate}

\begin{enumerate}
\item The reward model inherits this conflation, showing similar confidence levels on both task types (mean\_diff = 0.018).
\end{enumerate}

\begin{enumerate}
\item Hidden state analysis confirms models do not internally distinguish task types (separation = 0.024). A linear probe achieves only moderate accuracy (68-76\%), indicating weak task-type encoding.
\end{enumerate}

\begin{enumerate}
\item These RLHF models exhibit systematic calibration inversion patterns (silhouette = 0.6016).
\end{enumerate}

The gap between step 3 and step 4 remains: the causal link from bidirectional task features to calibration inversion cluster membership was not established with keyword-based detection.

\subsubsection{Cross-Model Consistency}

All three models (Llama-2-7B, Llama-2-13B, Mistral-7B) showed consistent patterns across all passing hypotheses. This cross-model consistency strengthens the finding that RLHF training creates consistent behavioral patterns regardless of base model architecture.

\subsubsection{Limitations}

\textbf{Keyword detection insufficient:} Simple keyword patterns (user\_belief, context\_dependent, hedged\_answer) captured only 2.3\% of tasks. Bidirectionality is a semantic property not reliably captured by lexical patterns.

\textbf{Open models only:} Only open RLHF models were tested. Closed models (GPT-4, Claude) may behave differently, but logprob access required for calibration analysis is not available for these models.

\textbf{Dataset composition:} TruthfulQA, MMLU moral\_scenarios, and Anthropic HH-RLHF may not contain sufficient natural bidirectional variation. These datasets were designed for accuracy evaluation, not bidirectionality analysis.

\textbf{Alternative explanations not tested:} Calibration inversion clusters may be driven by task difficulty, answer format, or topic rather than bidirectionality. These confounds were not systematically evaluated.

\subsubsection{Implications}

Calibration clustering offers a diagnostic tool for RLHF behavioral analysis. The finding that annotators provide no distinguishing signal between task types (conflation score = 0.999) suggests that annotation-level interventions may be more effective than post-hoc calibration adjustment.
